% !TeX program = xelatex
% UTF-8. Compile with XeLaTeX from this directory or the repository root.
% Transcription of Google Drive / Notebook / asthma; original filename order.
% previous-pic01.jpg: 20261016_054255746.jpg; Drive ID 1xskh1aUgnEctDY38hfHiCR-nQhBvYJFO
% previous-pic02.jpg: 20261016_054305290.jpg; Drive ID 13qnC1FU_AELQao3eqk6oUGtpkZU0iBNz
% previous-pic03.jpg: 20261016_054309081.jpg; Drive ID 1qEVbBRAp5zvDl0cy0eaCCvrevT5Vc-OR
% previous-pic04.jpg: 20261016_054315003.jpg; Drive ID 15x7ZmR1aKXMEMNJbTVQU3pNL6n29d1Lp
% previous-pic05.jpg: 20261016_054319779.jpg; Drive ID 13WfVwFf74y2IvqDFBFEfJ1u00jgzayN2
\documentclass[11pt,a4paper]{article}
\usepackage[margin=20mm]{geometry}
\usepackage{amsmath,amssymb}
\usepackage{graphicx}
\usepackage{enumitem}
\usepackage{xcolor}
\usepackage{xepersian}
\settextfont{Amiri}
\setlatintextfont{DejaVu Serif}
\setlength{\parindent}{0pt}
\setlength{\parskip}{4pt}
\setlength{\emergencystretch}{3em}
\setlist[itemize]{itemsep=3pt,topsep=3pt}
\newcommand{\amix}[1]{\lr{#1}}
\newcommand{\aread}[1]{\textcolor{red!65!black}{[خوانش نامطمئن: #1]}}
\newcommand{\ahead}[1]{%
  \section*{تصویر #1}
  \begingroup\small
  سربرگ: ۱۸۵؛ \amix{Subject:}\quad\amix{Date: / /} (تاریخ خالی).
  نشان چاپی پایین برگه: \amix{YASHA}؛ شمارهٔ دست‌نویس پایین: #1.
  \par\endgroup}

\begin{document}
\begin{center}
{\Large\bfseries آسم کودکان ـ ۱۸۵}\par
\amix{Nelson 22}
\end{center}

\ahead{۱}
\subsection*{عوامل خطر و تریگرها}
\begin{itemize}
  \item برونشیولیت در سال اول زندگی: یک \amix{RF} مهم برای آسم کودکی پس از ۵ سالگی.
  \item تنباکو در محیط، شدت آسم را بیشتر می‌کند.
  \item هوای سرد و خشک، هایپرونتیلاسیون طی ورزش، و \aread{بوهای شیمیایی}: تریگر.
\end{itemize}

\subsection*{اپیدمیولوژی}
\begin{itemize}
  \item ۱۱٪ کودکان در آمریکا؛ ۷۰٪ \aread{در بزرگسالی} هم آسم دارند.
  \item فقر و سیاه بودن: ریسک فاکتور.
  \item در ۸۰٪ موارد، قبل از ۶ سالگی شروع می‌شود.
  \item آلرژی در کودکان به همراه سرفه و/یا \amix{wheeze} راجعه، مهم‌ترین عامل برای ماندگار شدن آسم کودکی است.
\end{itemize}
\begingroup\small
در سطر اخیر، واژهٔ «عطسه» زیر \amix{wheeze} خط خورده و با همین واژهٔ انگلیسی جایگزین شده است.
\endgroup

\subsection*{ریسک فاکتورهای کودکی}
\begin{itemize}
  \item ماژور: آسم در والدین، اگزما، آلرژی به مواد استنشاقی.
  \item مینور: رینیت آلرژیک، ویز جدا از سرماخوردگی، ائوزینوفیلی \amix{\(>4\%\)} و آلرژی غذایی.
\end{itemize}

\subsection*{انواع آسم کودکان بر اساس سیر طبیعی علائم بالینی}
۱) ویز راجعه در ابتدای کودکی:
\begin{itemize}
  \item تریگر معمولاً عفونت وایرال.
  \item معمولاً قبل از مدرسه رفع می‌شود.
\end{itemize}

\clearpage
\ahead{۲}
۲) آسم مزمن:
\begin{itemize}
  \item با آلرژی ارتباط دارد.
  \item معمولاً تا بزرگسالی می‌ماند.
\end{itemize}
\begin{itemize}
  \item کودکان در سن مدرسه با آسم خفیف تا متوسط، معمولاً تا نوجوانی خوب می‌شوند.
  \item ۴۰٪ آن‌ها دوباره آسم \amix{intermittent} می‌شوند.
  \item \amix{ICS} برای آسم \amix{persistent}، ریسک ماندگاری آن تا بزرگسالی را تغییر نمی‌دهد.
\end{itemize}
\begingroup\small
سطر خط‌خوردهٔ زیر جملهٔ \amix{ICS}:
«حملات تشدید را \amix{\(\downarrow\)} و افت پیشروندهٔ عملکرد ریه را \aread{الغا می‌کند}».
\endgroup

\subsection*{دسته‌بندی آسم بر اساس شدت}
\begin{itemize}
  \item \amix{Intermittent}
  \item \amix{Persistent}:
  \begin{itemize}
    \item \amix{Mild}
    \item \amix{Moderate}
    \item \amix{Severe}
  \end{itemize}
\end{itemize}

\subsection*{دسته‌بندی بر اساس کنترل}
\begin{latin}
\begin{itemize}
  \item Well
  \item Not well
  \item Very poorly cont.
\end{itemize}
\end{latin}

\subsection*{بر اساس نیاز به دارو}
\begin{latin}
\begin{itemize}
  \item Easy to control
  \item Difficult to control
  \item Exacerbators
  \item Refractory
\end{itemize}
\end{latin}

\clearpage
\ahead{۳}
\subsection*{تظاهرات بالینی و تشخیص}
\begin{itemize}
  \item اپیزودهای مکرر سرفهٔ خشک و ویز بازدمی.
  \item در کودکان بزرگ‌تر:
  \begin{itemize}
    \item تنگی نفس \amix{(SoB)}.
    \item گرفتگی سینه \amix{(Chest Congestion)}.
    \item درد قفسهٔ سینهٔ \amix{non focal}.
  \end{itemize}
  \item علائم تنفسی در شب بدتر می‌شوند.
  \item در روز، به دنبال فعالیت تشدید می‌شود.
  \item عدم پاسخ به برونکودیلاتور و \amix{CS}، ردکنندهٔ آسم است.
  \item تریگرهای علائم: فعالیت فیزیکی، هایپرونتیلاسیون (خنده)، هوای سرد و خشک، عفونت‌های تنفسی.
  \item بیماران آسمی در زمان بررسی معمولاً بدون علامت هستند.
  \item رفع سریع علائم (ظرف \amix{10 min}) به دنبال \amix{SABA} استنشاقی به نفع آسم است؛ آلبوترول.
\end{itemize}

\subsection*{سمع}
\begin{itemize}
  \item در سمع ممکن است در ناحیهٔ \amix{Right Lower Posterior Field}، کاهش صدا رخ دهد.
  \item رونکای و کراکل (رال): به دلیل مخاط زیاد، آتلکتازی نقطه‌ای.
  \item در اگزا‌سربیشن: ویز هم دم هم بازدم.
  \item خیلی شدید: \amix{Silent chest}.
\end{itemize}

\clearpage
\ahead{۴}
\begin{itemize}
  \item تست‌های آزمایشگاهی برای بررسی آلرژی مفید هستند ولی برای تشخیص آسم جایگاهی ندارند.
  \item \amix{PFT} تأیید تشخیص آسم است.
\end{itemize}

\begin{latin}
\subsection*{PFT in Asthma}
\begin{itemize}
  \item Spirometry \(\longrightarrow\) Airflow Limitation
  \begin{itemize}
    \item Low \(\mathrm{FEV}_1\) (Many patients with asthma have N/L \(\mathrm{FEV}_1\)).
    \item \(\mathrm{FEV}_1/\mathrm{FVC}\) ratio below the lower limit.
  \end{itemize}
  \item Bronchodilator response
  \begin{itemize}
    \item Reversible if: \(\mathrm{FEV}_1\uparrow > 9\text{--}12\%\) or predicted \(\mathrm{FEV}_1\uparrow >10\%\).
  \end{itemize}
  \item Exercise challenge \(\longrightarrow\) \(\mathrm{FEV}_1\downarrow \geq 15\%\).
  \item Daily Peak Expiratory Flow (PEF) or \(\mathrm{FEV}_1\) monitoring
  \begin{itemize}
    \item \(\geq 20\%\) Variation.
  \end{itemize}
\end{itemize}
\end{latin}
\begingroup\small
در سطر \amix{predicted}، حروف شاخص در دست‌نویس به شکل \amix{FVE\(_1\)} آمده‌اند؛
در سطر \amix{Exercise challenge} نیز یک علامت پیش از \amix{\(\geq\)} خط خورده است.
\endgroup

\subsection*{کاربرد \amix{PFT}}
\begin{itemize}
  \item \amix{PFT}: تشخیص و ارزیابی درمان.
  \item در کودکان ۴–۵ و بالاتر قابل اجرا است.
  \item ۵۰٪ کودکان با آسم خفیف تا متوسط، اسپیرومتری \amix{N/L} دارند.
  \item چالش فعالیت: دویدن برای ۶–۸ دقیقه.
\end{itemize}

\clearpage
\ahead{۵}
\subsection*{\amix{EIB}}
\begin{itemize}
  \item \amix{EIB}: برونکواسپاسم در زمینهٔ فعالیت.
  \item از دقیقهٔ ۵ شروع؛ پیک دقیقهٔ ۱۵؛ رفع ظرف ۳۰–۶۰.
\end{itemize}

\subsection*{رادیوگرافی در آسم}
\begin{itemize}
  \item معمولاً \amix{N/L} است.
  \item شواهد هایپرونتیلاسیون، هایپرانفلیشن:
  \begin{itemize}
    \item \amix{Flat} شدن دیافراگم.
    \item پری برونکیال \amix{thickening}.
  \end{itemize}
  \item کاربرد اصلی بررسی \amix{DDX} است.
  \item برونشکتازی در \amix{CT} بهتر دیده می‌شود.
\end{itemize}

\clearpage
% BEGIN added six-page transcription: Drive Notebook/asthma, 20261018.
% pic01.jpg: 20261018_195633193.jpg; Drive ID 1b-eYNb_iEzoT_HwyYhdUkNcptmp7bZiQ
% pic02.jpg: 20261018_195638511.jpg; Drive ID 1N9Z_vxH4rvAEOJcL8-29tuzzxlPjpsj9
% pic03.jpg: 20261018_195646775.jpg; Drive ID 1HwmZaqLRZgPV0axdb5QikuzIKviL37Oa
% pic04.jpg: 20261018_195651668.jpg; Drive ID 1y5-HjP6t9Q5DwMqjxnK00UQqAqP5T_3A
% pic05.jpg: 20261018_195659061.jpg; Drive ID 10WW-xrkgWPOWC9pM3TAWubE0-dkdb0rB
% pic06.jpg: 20261018_195703686.jpg; Drive ID 1IxntzrvM-XYEn_tGQpS57Cl1ZLxRomU_
\section*{آسم ـ ۷۸؛ مجموعهٔ افزوده}
\begingroup\small
سربرگ مشترک هر شش تصویر: \amix{Subject:} «۷۸ ـ آسم»؛
\amix{Year:}\quad\amix{Month:}\quad\amix{Date:} (هر سه خالی).
شماره‌های چاپی سطرها: ۱ تا ۹؛ شماره‌های دست‌نویس پایین صفحات: ۱ تا ۶.
خوانش‌های نامطمئن و نوشته‌های بریده‌شده در حاشیه با کروشه مشخص شده‌اند.
\par\endgroup

\subsection*{تصویر افزودهٔ ۱ ـ کلیات و تظاهرات بالینی}
\begin{itemize}
  \item ازدیاد حساسیت راه‌های هوایی نسبت به: آلرژن‌ها، محرک‌ها، عفونت‌های ویروسی، فعالیت.
  \item التهاب مزمن طی آسم: ریمدلینگ راه‌های هوایی.
  \item آسم در کودکان: \amix{\(F<M\)}؛ در بزرگسالان: \amix{\(M<F\)}.
  \item تظاهرات بالینی: سرفه، ویزینگ، \amix{SoB}، عدم تحمل ورزش؛
  پس از «ورزش» یک واژهٔ خط‌خورده دیده می‌شود: \aread{سرفه}.
  می‌تواند علائم شبانه داشته باشد.
  \item رینوسینوزیت، \amix{GERD}، آسپرین: آسم را تشدید می‌کند.
  \item طی اپیزودهای آسم: تاکی‌پنه، تاکی‌کاردی، سرفه، ویزینگ، بازدم طولانی.
  \item حمله‌های شدید: سیانوز، رترکشن، آژیتیشن، عدم توانایی صحبت کردن،
  پوزیشن نشستن تریپاد، دیافورز، پالس پارادوکس.
  \item پالس پارادوکس: کاهش بیش از \amix{\(15\,\mathrm{mmHg}\)} فشار خون طی دم.
  \item اسپیرومتری در کودکان بالای ۵ سال قابل انجام است.
\end{itemize}

\subsection*{\amix{DDX} آسم}
\begin{itemize}
  \item بیماری‌هایی که انسداد راه‌های هوایی بزرگ می‌دهند:
  \begin{itemize}
    \item جسم خارجی.
    \item اثر فشاری تومور یا لنف نود.
    \item وب لارنژیال، رینگ واسکولار.
    \item حرکات پارادوکسیکال تار صوتی.
  \end{itemize}
  \item بیماری‌هایی که انسداد راه‌های هوایی کوچک می‌دهند:
  \begin{itemize}
    \item \amix{CF}.
    \item دیسپلازی برونکوپولمونری (کرانیک لان دیزیز آف پره‌مچوریتی).
    \item برونشیولیت‌ها.
  \end{itemize}
  \item \amix{GERD}، \amix{PND}، رینوسینوزیت آلرژیک، ادم ریه، \amix{HF}.
\end{itemize}

\clearpage
\subsection*{تصویر افزودهٔ ۲ ـ علل سرفه و بررسی‌ها}
\subsubsection*{علل سرفه در ماه‌های اول تولد}
\begin{itemize}
  \item \amix{C: CF}
  \item \amix{R: RTI}
  \item \amix{A}: آسپیریشن.
  \item \amix{D}: دیس‌فانکشن سیلیا.
  \item \amix{L: Lung and airway malformation}: استنوز تراکئال، لارنگوتراکئومالاسی، وب.
  \item \amix{E: Edema}: \amix{HF}، بیماری قلبی مادرزادی.
\end{itemize}

\subsubsection*{\amix{FeNO}}
\begin{itemize}
  \item \amix{FeNO} (نیتریک اکساید بازدمی): نشان‌دهندهٔ التهاب ائوزینوفیلیک مجاری هوایی.
  \item نشان‌دهنده و پیش‌بینی‌کنندهٔ میزان پاسخ به \amix{CS}.
  \item برای اصلاح دوز دارو استفاده می‌شود.
\end{itemize}

\subsubsection*{تست آلرژی پوستی}
\begin{itemize}
  \item در همهٔ کودکان با آسم ماندگار باید انجام شود.
  \item در اگزرسبیشن نباید انجام شود.
\end{itemize}

\subsubsection*{\amix{CXR}}
\begin{itemize}
  \item در اولین اپیزود آسم یا دومین اپیزود سرفه یا ویزینگ، به عنوان تشخیص و برای بررسی اختلالات آناتومیک.
  \item تکرار \amix{CXR} در اپیزودهای بعدی لازم نیست مگر تب وجود داشته باشد؛ برای بررسی پنومونی.
\end{itemize}

\subsubsection*{آلرژیک برونکوپولمونری آسپرژیلوزیس}
\begin{itemize}
  \item واکنش ازدیاد حساسیت به \amix{Ag} قارچ آسپرژیلوس.
  \item در بیمارانی که مبتلا به آسم هستند و \amix{CS} مصرف می‌کنند و بیماران \amix{CF} دیده می‌شود.
\end{itemize}

\begin{itemize}
  \item کودک مبتلا به آسم: از دود سیگار، چوب و فرد مبتلا به عفونت ویروسی دوری کند.
  \item واکسیناسیون آنفولانزا.
\end{itemize}

\clearpage
\subsection*{تصویر افزودهٔ ۳ ـ داروهای طولانی‌اثر}
\subsubsection*{\amix{ICS} طولانی‌اثر}
\begin{itemize}
  \item مؤثرترین داروها.
  \item قد کودکان باید بررسی شود؛ در سال اول \amix{\(1\,\mathrm{cm}\)} کمتر رشد می‌کند.
  یادداشت حاشیه: «برگشت‌پذیر \aread{ادامه در لبهٔ تصویر بریده شده}».
  \item بر متابولیسم اثری ندارد.
  \item عارضهٔ چشمی ندارد.
  \item باید کمترین دوز ممکن دریافت شود.
\end{itemize}

\subsubsection*{لوکوترین مودیفایرها}
\begin{itemize}
  \item لوکوترین‌ها مدیاتور التهاب، عامل اسپاسم عضلات صاف تنفسی هستند.
  \item دو کلاس دارند:
  \begin{itemize}
    \item آنتاگونیست رسپتور لوکوترین: مونته‌لوکاست، زفیرلوکاست.
    \item مهارکنندهٔ سنتز لوکوترین: \amix{Zileuton}.
  \end{itemize}
  \item زفیرلوکاست: ۲ بار در روز؛ کودکان بالای ۵ سال.
  \item مونته‌لوکاست: ۱ بار در روز، شب‌ها.
  \item در آسم خفیف مؤثر است و کنترل فعالیت را بهتر می‌کند.
\end{itemize}

\subsubsection*{\amix{LABA} (آگونیست \amix{\(\beta_2\)} طولانی‌اثر)}
\begin{itemize}
  \item اثر ضدالتهاب ندارند.
  \item علائم را برای ۱۲–۲۴ ساعت ریلیف می‌کند.
  \item در درمان تک‌دارویی کاربرد ندارد: \amix{LABA + ICS}.
  \item فورموترول: بالای ۵ سال؛ شروع اثر ۱۵ دقیقه.
  \item آرفورموترول: شروع اثر ۱۵ دقیقه؛ ۲ بار در روز.
  \item سالمترول: بالای ۴ سال؛ شروع اثر ۳۰ دقیقه؛ ۲ بار در روز.
  \item ویلانترول: بسیار طولانی‌اثر؛ بالای ۱۸ سال؛ ۱ بار در روز.
\end{itemize}

\subsubsection*{بیولوژیک‌ها}
\begin{itemize}
  \item \amix{Omalizumab}: \amix{anti IgE Ab}؛ مهار اتصال \amix{IgE} به ماست سل و بازوفیل و
  \aread{ادامهٔ عبارت در حاشیهٔ چپ تصویر بریده شده است}.
  تزریق \amix{SC} هر ۲–۴ هفته.
  \item \amix{Mepolizumab}: آسم شدید با فنوتیپ ائوزینوفیلیک؛ کاهش تولید و بقا ائوزینوفیل‌ها.
\end{itemize}

\clearpage
\subsection*{تصویر افزودهٔ ۴ ـ درمان‌ها و دسته‌بندی}
\subsubsection*{درمان‌های کوتاه‌اثر}
\begin{itemize}
  \item \amix{SABA}: آلبوترول، لوالبوترول؛ شروع اثر ۵–۱۰ دقیقه؛ طول اثر ۴–۶ ساعت.
  \item داروهای آنتی‌کولینرژیک:
  \begin{itemize}
    \item ایپراتروپیوم بروماید: برونکودیلاتور کوتاه‌اثر.
    \item تولید موکوس را کاهش می‌دهد.
    \item گیرنده‌های سرفه را سرکوب می‌کند.
    \item به کنار \amix{SABA} در حملات آسم اضافه می‌شود.
  \end{itemize}
  \item تیوتروپیوم: آنتی‌کولینرژیک طولانی‌اثر؛ آسم متوسط تا شدید در کودکان بالای ۶ سال.
  \item \amix{OCS}: برای حملهٔ آسم به مدت ۱–۱۰ روز استفاده می‌شود.
\end{itemize}
\begin{latin}
\begin{itemize}
  \item \(1\text{--}2\,\mathrm{mg/kg/day}\) of prednisone for 5 days.
  \item \(0.6\,\mathrm{mg/kg/day}\) of dexamethasone for 1--2 days.
\end{itemize}
\end{latin}

\subsubsection*{دسته‌بندی سنی برای درمان آسم}
\begin{itemize}
  \item ۰–۴ سال.
  \item ۵–۱۱ سال.
  \item \amix{\(\geq12\)} سال.
\end{itemize}

\begin{itemize}
  \item \amix{SABA} برای همه باید تجویز شود؛ استفادهٔ بیش از ۲ بار در هفته = درمان ناکافی.
  \item آسم اینترمیتنت: علائم کمتر از ۲ بار در هفته؛ \amix{SABA PRN}.
  \item آسم پرسیستنت: علائم روزانه حداقل ۲ بار در هفته؛ یا بیدار شدن از خواب حداقل ۲ بار در ماه.
  اندیکاسیون درمان ضدالتهابی روزانه.
\end{itemize}

\clearpage
\subsection*{تصویر افزودهٔ ۵ ـ مراحل درمان و اگزرسبیشن}
\subsubsection*{مراحل درمان در ۰–۴ سال}
\begin{latin}
\begin{enumerate}[label=Step \arabic*:]
  \item SABA PRN.
  \item Low dose ICS.
  \item Medium dose ICS.
  \item Medium ICS + LABA or Monte.
  \item High ICS + LABA or Monte.
  \item High ICS + LABA or Montel + OCS.
\end{enumerate}
\end{latin}

\subsubsection*{مراحل درمان در ۵–۱۱ سال}
\begin{latin}
\begin{enumerate}[label=Step \arabic*:]
  \item SABA PRN.
  \item Low dose ICS.
  \item Low dose ICS + LABA, LTRA, Theo.
  \item Medium ICS + LABA.
  \item High ICS + LABA.
  \item High ICS + LABA + OCS.
\end{enumerate}
\end{latin}

\begin{itemize}
  \item برای پایین آمدن، حداقل ۳ ماه کنترل باشد.
  \item وقتی آسم کنترل شد: ویزیت هر ۱–۶ ماه؛ قبل از کنترل، هر ۴–۶ هفته.
\end{itemize}

\subsubsection*{استاتوس آسماتیکوس و اگزرسبیشن آسم}
\begin{itemize}
  \item استاتوس آسماتیکوس: \amix{ABG} می‌گیریم.
  \item اگر اگزرسبیشن آسم که به درمان سرپایی جواب نمی‌دهد: بستری نیاز است.
  \item نشانه‌های اگزرسبیشن آسم: دیسترس تنفسی، دیس‌پنه، ویزینگ، سرفه.
\end{itemize}
\begin{latin}
Dicrease in spirometry or peak expiratory flow rate (PEFR).
\end{latin}
\begin{itemize}
  \item در موارد شدید اگزرسبیشن آسم: پالس اکسیمتری و \amix{ABG}.
  \item در حضور تاکی‌پنه، \amix{\(\mathrm{PCO}_2\)} طبیعی
  (\amix{\(35\text{--}45\,\mathrm{mmHg}\)}):
  \aread{ارست تنفسی نزدیک است؛ انتهای عبارت در لبهٔ چپ تصویر قرار دارد}.
  \item درمان اگزرسبیشن آسم:
  \amix{\(\mathrm{O}_2\)} تراپی + \amix{SABA} + کورتیکواستروئید خوراکی یا
  \aread{مسیر دوم تجویز در لبهٔ چپ تصویر بریده شده است}؛
  ایپراتروپیوم هم می‌دهیم.
  \item اگر جواب نداد: منیزیم سولفات تزریقی:
  \amix{\(25\text{--}75\,\mathrm{mg/kg}\) (max \(2.0\,\mathrm{g}\)) IV over 20 min}.
  \item اگر اگزرسبیشن آسم همراه با آنافیلاکسی: اپی‌نفرین \amix{IM} یا تربوتالین \amix{SC}.
\end{itemize}

\clearpage
\subsection*{تصویر افزودهٔ ۶ ـ پیش‌آگهی و ادامهٔ مراحل درمان}
\begin{itemize}
  \item آتوپی قوی‌ترین پیش‌بینی‌کنندهٔ تبدیل شدن به آسم دائمی در کودکان زیر ۳ سال است.
\end{itemize}

\subsubsection*{مراحل درمان در \amix{\(\geq12\)} سال}
\begin{latin}
\begin{enumerate}[label=Step \arabic*:]
  \item SABA PRN.
  \item Low dose ICS.
  \item Low dose ICS + LABA or Medium ICS.
  \item Medium ICS + LABA.
  \item High ICS + LABA and Omalizumab.
  \item High ICS + LABA + OCS and Omalicumab.
\end{enumerate}
\end{latin}
یادداشت زیر \amix{Omalizumab} در مرحلهٔ ۵: «اگر آلرژی».
املای \amix{Omalicumab} در مرحلهٔ ۶ مطابق دست‌نویس نگه داشته شده است.

\begin{itemize}
  \item \amix{Peak flow monitoring} در کودکان بالای ۵ سال.
\end{itemize}

\begin{latin}
\subsubsection*{Modified asthma predictive index for children}
\begin{itemize}
  \item At least 4 wheezing episodes + One major or Two minor criteria.
  \item M: Parental asthma, Eczema, Inhalant Allergen sensitization.
\end{itemize}
\end{latin}
\begin{itemize}
  \item \amix{m}: ائوزینوفیلی \amix{\(\geq4\%\)}، ویزینگ، رینیت آلرژیک،
  \amix{Food Allergen sensitization}.
\end{itemize}
% END added six-page transcription.

\clearpage
\section*{تصاویر اصلی}
\IfFileExists{images/previous-pic01.jpg}{\def\aimagebase{images/}}{%
  \def\aimagebase{pediatrics/ped-pulmonology/asthma/images/}}
\newcommand{\aoriginal}[2]{%
  \begin{center}
    \textbf{تصویر #1}\par
    \includegraphics[width=\linewidth,height=0.88\textheight,keepaspectratio]{\aimagebase #2}
  \end{center}}
\aoriginal{۱}{previous-pic01.jpg}
\clearpage
\aoriginal{۲}{previous-pic02.jpg}
\clearpage
\aoriginal{۳}{previous-pic03.jpg}
\clearpage
\aoriginal{۴}{previous-pic04.jpg}
\clearpage
\aoriginal{۵}{previous-pic05.jpg}
\clearpage
\section*{تصاویر اصلی مجموعهٔ افزوده}
\aoriginal{۱}{pic01.jpg}
\clearpage
\aoriginal{۲}{pic02.jpg}
\clearpage
\aoriginal{۳}{pic03.jpg}
\clearpage
\aoriginal{۴}{pic04.jpg}
\clearpage
\aoriginal{۵}{pic05.jpg}
\clearpage
\aoriginal{۶}{pic06.jpg}
\end{document}
